%!TEX root = ../dokumentation.tex

\chapter{Conclusion}
\label{ch:conclusion}

\section{Summary of Findings}

This thesis compared classical and agent-based approaches to knowledge retrieval across a heterogeneous enterprise documentation corpus through four distinct system configurations, each building on the limitations identified in its predecessor. Both pipelines were implemented from scratch, shared the same embedding model, chunking strategy, reranker, and LLM, and were evaluated over 90 queries spanning four types, namely factual, synthesis, ambiguous, and unanswerable, against three simulated enterprise source types.

The initial two-system comparison between the original RAG pipeline and the single-agent ReAct system confirmed RAG's advantage on every measurable dimension. The RAG pipeline achieved a context precision of 0.546 on factual queries and 0.730 on synthesis queries, against the agent's 0.338 and 0.315. Answer relevancy followed the same pattern: RAG reached 0.792 on factual and 0.679 on synthesis queries; the agent reached 0.663 and 0.547 respectively, with 41\,\% of queries producing empty answers due to the ReAct iteration ceiling. The latency and cost differences were equally stark: the agent required 6.5$\times$ longer per query and consumed 11.2$\times$ more tokens. The agent's only structural advantage, namely auditable source routing, achieved a 73.3\,\% accuracy that did not translate into quality gains sufficient to offset these costs.

The multi-agent coordinator architecture, implemented as a direct response to the single agent's failure modes, eliminated every failure the single agent exhibited. By replacing the sequential ReAct cycle with a two-node LangGraph graph comprising a retrieval node that queries all three sources sequentially followed by a single synthesis call over a bounded 15-chunk context, the multi-agent achieved a 0\,\% failure rate and produced the highest answer relevancy of any system evaluated: 0.818 on factual queries and 0.770 on synthesis queries, compared to original RAG's 0.792 and 0.679 respectively. Latency fell to 4.3 seconds mean and cost to \$2.91 per 90 queries, representing 2.5$\times$ RAG's cost rather than 11.2$\times$. Unanswerable faithfulness reached 0.839, the highest of any system, reflecting the bounded context's resistance to fabrication when no relevant material is retrieved.

The improved RAG pipeline, which applied HyDE and MMR to the original pipeline, demonstrated that significant quality gains are available within the classical architecture without structural change. Synthesis answer relevancy improved from 0.679 to 0.742, a gain of 9\,\%, at a cost increase of \$0.02 per 90 queries. Latency increased from 2.5 to 4.3 seconds due to the extra hypothesis generation call. BERTScore and ROUGE-L were insensitive to these improvements, confirming that reference-based lexical metrics underestimate the practical benefit of advanced retrieval techniques in open-domain settings.

The complexity rubric assigns mean scores of 4.0 for both RAG variants, 3.2 for the multi-agent coordinator, and 2.2 for the single-agent system. The multi-agent is more complex to build and maintain than either RAG variant, but substantially more transparent and debuggable than the single-agent design it supersedes. Adding a new data source requires two steps for RAG and three for the multi-agent, compared to five for the single agent, where routing behaviour must be revalidated across all existing tools whenever the toolset changes.


\section{Practical Recommendation}

The four-system evaluation produces a differentiated practical recommendation that depends on the deployment context and its primary constraints.

\textit{The original RAG pipeline} remains the appropriate starting point for most enterprise deployments. It is the fastest to implement, the cheapest to run (\$1.18 per 90 queries, 2.5-second mean latency), and the simplest to debug and extend. For query mixes dominated by factual lookups against a coherent knowledge domain, no other configuration evaluated here offers a better cost-quality ratio.

\textit{The improved RAG pipeline} (with HyDE and MMR) is the right incremental upgrade when synthesis quality is a stated requirement and the cost and infrastructure must remain unchanged. The 9\,\% synthesis answer relevancy gain at negligible additional cost, achieved without new services, dependencies, or architectural changes, makes this the highest-return improvement available to a team operating an existing RAG deployment. A team already running the original RAG pipeline can apply HyDE and MMR by modifying three files (\texttt{config.py}, \texttt{retrieve.py}, \texttt{generate.py}), without touching the vector store, embedding model, or evaluation infrastructure.

\textit{The multi-agent coordinator} is the preferred architecture when maximum answer quality across all query types is the priority and a 2.5$\times$ cost premium over RAG is acceptable. It delivers the best answer relevancy on both factual (0.818) and synthesis (0.770) queries, handles unanswerable queries with the highest faithfulness of any system (0.839), produces no empty-answer failures, and provides full source attribution via the structured \texttt{agent\_trace} log. Its main operational constraints are the requirement for three maintained retrieval indexes and the inability to selectively target a subset of sources, as it always queries all three, which may introduce noise for queries clearly targeting a single domain.

\textit{The single-agent ReAct system} as evaluated here should not be deployed for general knowledge retrieval. It offers no quality advantage over any alternative, its 41\,\% failure rate makes aggregate quality metrics unreliable, and its cost and latency are prohibitive for production use. The failure modes documented in this thesis, comprising context accumulation exceeding the iteration ceiling, reasoning tokens exhausting the answer budget, and the absence of an exit path for unanswerable queries, are not incidental to a particular implementation; they are structural consequences of the sequential ReAct cycle combined with a reasoning-heavy model. Any future deployment of a sequential agentic retrieval system must address all three failure modes explicitly before evaluation.

In practice, a graduated adoption path reflects the findings most directly: deploy the improved RAG pipeline as the production baseline, monitor synthesis failure cases through structured logging, and escalate to the multi-agent coordinator for query types or user segments where answer completeness consistently underperforms. This path minimises both initial implementation cost and operational risk while preserving the option to increase quality investment where it demonstrably matters.


\section{Future Work}

Three directions offer the highest return for extending this evaluation.

\textbf{Real enterprise data validation.} The proxy corpus approximates enterprise knowledge systems but cannot reproduce their full complexity: proprietary terminology, access-control fragmentation, inconsistent document quality, and mixed languages. Replicating this evaluation on a live enterprise corpus, with appropriate data governance controls, would test whether the quality and latency relationships observed here hold under real-world noise conditions. In particular, the multi-agent's always-all-sources design may behave differently when one source type contains terminology that interferes with synthesis queries targeting another; this effect is not observable in the thematically coherent cloud-native corpus used here.

\textbf{Hybrid routing for the multi-agent architecture.} The multi-agent's 100\,\% routing coverage is an architectural property rather than a capability, achieved by always querying all sources rather than by predicting which source is relevant. This introduces unnecessary latency and potential context dilution on queries with a clear single-source target. A lightweight query classifier trained to predict source relevance, routing clear single-source queries to one retriever while directing ambiguous or synthesis queries to all sources, would combine the multi-agent's reliability with the efficiency of targeted retrieval. The structured logs and routing ground-truth labels produced in this evaluation provide a natural training and evaluation set for such a classifier.

\textbf{Cost and latency optimisation for HyDE.} The improved RAG pipeline's latency increase from 2.5 to 4.3 seconds is attributable entirely to the HyDE hypothesis generation call, which uses the same \texttt{gpt-oss-120b} reasoning model as the answer generation step. Using a smaller, faster model for hypothesis generation, one capable of producing a fluent 2--3 sentence passage in a target domain without chain-of-thought reasoning, would reduce the latency overhead of HyDE substantially without degrading hypothesis quality. The evaluation framework developed here, which captures per-stage latency and token counts in structured logs, provides the instrumentation needed to measure this trade-off precisely without re-running the full 90-query evaluation after each model substitution.
